Una domanda relativa a un'emergenza naturale raramente trova risposta in una
sola fonte. Sapere se si è verificato un terremoto, verificare quali acquisizioni
satellitari coprano la zona e controllare se sia stata attivata una produzione
cartografica sono operazioni diverse. I dati possono riguardare lo stesso luogo
e lo stesso intervallo temporale, ma descrivono fenomeni distinti e hanno tempi
di aggiornamento, formati e livelli di affidabilità propri.

Si consideri, ad esempio, una richiesta sui terremoti avvenuti in Italia
nell'ultima settimana e sulla relativa copertura satellitare. La prima parte
richiede un elenco di eventi e dei rispettivi livelli di allerta; la seconda
riguarda la disponibilità di osservazioni della Terra. Un'eventuale attivazione
di un servizio di emergenza aggiunge poi informazioni su aree di interesse e
prodotti cartografici. Riunire questi elementi in una sola risposta è utile
soltanto se la loro provenienza rimane riconoscibile.

Questa tesi affronta il problema realizzando \emph{Copernicus Data Workspace}:
uno spazio di lavoro conversazionale che coordina fonti geospaziali indipendenti.
L'utente formula la richiesta in linguaggio naturale; il sistema ne ricava
luogo, periodo e obiettivo, sceglie le fonti pertinenti e restituisce testo,
dati strutturati e visualizzazioni. Il modello linguistico fa da interprete:
comprende la domanda e formula la risposta, mentre gli eventi e i valori
mostrati nei grafici provengono dalle fonti consultate.

\section{Posizionamento rispetto ai lavori correlati}

Il progetto si colloca nella linea degli agenti basati su LLM che usano
strumenti esterni. ReAct studia l'alternanza tra ragionamento e azione;
Toolformer affronta l'apprendimento dell'uso delle API
\cite{yaoReAct,schickToolformer}. Il function calling documenta invece un
meccanismo applicativo per collegare il modello a funzioni eseguibili
\cite{openaiFunctionCalling}.

Nel dominio dell'osservazione della Terra, Earth Copilot nasce dalla
collaborazione NASA--Microsoft per rendere interrogabili dati geospaziali
attraverso domande in linguaggio naturale \cite{earthCopilot2024}; la
documentazione pubblica ne descrive finalità e integrazione con l'accesso ai
dati NASA. Il CDSE STAC Browser offre invece navigazione e filtri sui metadati
del catalogo, con accesso alle proprietà e agli asset dei prodotti
\cite{cdseStac}. È un termine di confronto per la consultazione dei cataloghi,
mentre il workspace qui proposto aggiunge interpretazione linguistica,
coordinamento di fonti eterogenee e continuità conversazionale.

Il posizionamento del lavoro riguarda l'integrazione di questi bisogni:
Discovery di risorse catalogate, ricerca STAC live e consultazione di eventi
e attivazioni. MCP \cite{mcpSpec} delimita l'accesso a GDACS e CEMS;
il catalogo a grafo conserva le gerarchie delle risorse e il ranking applica
regole esplicite. Il contributo è un'architettura applicativa, costruita su
tecniche di retrieval esistenti, insieme al relativo protocollo di verifica.

\section{Il problema affrontato}

Le applicazioni conversazionali rendono semplice l'accesso a servizi che, presi
singolarmente, richiederebbero la conoscenza di endpoint, parametri e formati di
risposta. La complessità dei dati sottostanti, però, rimane. Nel percorso
operativo sono coinvolte tre famiglie di informazioni, alle quali si affianca
un catalogo trasversale di metadati.

GDACS pubblica eventi e allerte relativi a disastri naturali. Il catalogo del
Copernicus Data Space Ecosystem descrive invece le acquisizioni satellitari
disponibili per una certa area e un certo periodo. Copernicus Emergency
Management Service documenta attivazioni e prodotti cartografici predisposti a
supporto della gestione delle emergenze. Queste fonti possono essere consultate
nello stesso turno e ciascuna conserva il proprio significato: il numero di
acquisizioni satellitari, per esempio, misura la disponibilità di osservazioni,
mentre una prima stima della gravità di un evento viene dai livelli di
allerta GDACS.

Alla distinzione semantica si aggiunge quella tecnica. GDACS distribuisce anche
feed XML, il catalogo satellitare espone un'interfaccia STAC e CEMS adotta
strutture specifiche per le proprie attivazioni. I servizi possono restituire
risposte vuote, diventare temporaneamente indisponibili o richiedere trattamenti
diversi degli errori. Collocare questi dettagli nei prompt renderebbe difficile
capire se un'informazione sia stata recuperata oppure generata dal modello.

Anche l'interpretazione linguistica deve essere verificata. Un LLM sa trattare
espressioni come ``la settimana scorsa'', riconoscere un codice CEMS e comporre
un testo leggibile, ma la sua conoscenza può essere incompleta o datata: una
frase plausibile va riscontrata nella fonte consultata. Nel progetto
interpretazione, accesso ai servizi e calcoli sono perciò passaggi separati.

I turni successivi richiedono inoltre di conservare il contesto. Dopo una tabella
o un grafico, l'utente può chiedere di mostrare una sola piattaforma, cambiare
l'aggregazione o aggiornare una sorgente. Il significato di questi comandi
dipende dai risultati precedenti; lo stato è quindi associato alla
conversazione, mentre i dati remoti restano fuori dalla cronologia testuale.

La verificabilità riguarda sia il testo sia le visualizzazioni. Link,
identificativi, timestamp e criteri di aggregazione fanno parte del risultato.
Un grafico deve indicare i record e i campi dai quali deriva; il modello può
proporre la vista, mentre conteggi e misure vengono calcolati dal software.

\section{Obiettivo del lavoro}

L'obiettivo è un prototipo che trasformi una richiesta geografica in un percorso
di consultazione controllato. Il sistema deve poter interrogare più fonti quando
la domanda lo richiede, continuare a produrre una risposta anche quando una
sorgente secondaria è indisponibile e mantenere il contesto necessario per i
turni successivi. I risultati devono restare associati alla fonte dalla quale
provengono.

La tesi impiega modelli linguistici esistenti e studia l'architettura che li
circonda: ammissione della domanda nel dominio, risoluzione dei parametri,
scelta delle sorgenti e restituzione dei risultati nel workspace. Le ambiguità
linguistiche sono affidate al modello; schemi, filtri, calcoli e persistenza
rimangono sotto il controllo del codice.

Il contributo dichiarato è \textbf{progettuale e implementativo}: un prototipo
con responsabilità specializzate, percorsi Discovery/operativo controllati e
artefatti verificabili. Il titolo parla di approccio multiagentico perché il
workflow coordina componenti decisionali basati su LLM con responsabilità
specializzate, affiancati da nodi deterministici di recupero e trasformazione.

Il percorso \textbf{operativo} combina la ricerca web con tre fonti strutturate.
GDACS fornisce eventi e allerte; il catalogo STAC di Copernicus descrive
acquisizioni satellitari e relativa copertura; CEMS permette di consultare
attivazioni di Rapid Mapping e Risk and Recovery Mapping. Il contenuto della
domanda determina quali percorsi siano necessari e più rami possono essere
attivati nello stesso turno.

Il percorso di \textbf{Discovery} affronta una domanda complementare: quali
risorse Copernicus permettono di studiare una certa variabile, area o periodo?
Tre harvester raccolgono metadati Marine, CDSE/STAC e Drought in NebulaGraph.
Il workflow consulta questo grafo per individuare dataset, variabili,
collezioni e layer pertinenti. La Discovery indica quali risorse esistono, per
esempio un dataset di temperatura, e come raggiungerle; la ricerca delle
singole acquisizioni resta affidata al servizio STAC remoto. Sono ammessi anche
percorsi \textbf{combinati}, in cui Discovery e analisi operativa convivono
nella stessa sessione mantenendo distinti metadati di catalogo e scene
satellitari live.

Per GDACS e CEMS è stato adottato il Model Context Protocol. I dettagli dei due
servizi vengono incapsulati in server separati, ai quali il backend accede
attraverso un client comune. Il catalogo STAC live rimane integrato direttamente
nell'applicazione. La presenza dei due approcci nello stesso progetto permette
di confrontarne gli effetti sull'organizzazione del codice.

Il punto di ingresso è una dashboard web autenticata con Keycloak: home in
italiano con selettore IT/EN, formulazione della domanda, suggerimenti e
ripresa delle indagini recenti. La risposta operativa può includere una
dashboard delle acquisizioni Copernicus. Grafici, indicatori, tabella e mappa
derivano dai record STAC normalizzati; GDACS, CEMS e ricerca web restano nella
parte testuale o in sezioni dedicate, perché le rispettive quantità hanno
natura diversa. Durante la conversazione l'utente può applicare filtri, cambiare
vista o chiedere una nuova aggregazione. Ogni modifica produce una revisione
dell'artefatto e i valori sono sempre ricalcolati sul dataset conservato.

La Discovery restituisce invece un artefatto di risorse catalogate con
identificativi, descrizioni e collegamenti. Sono quindi distinte la vista dei
metadati e quella delle acquisizioni. A entrambe si affiancano un'infrastruttura
di osservabilità, per ricostruire tempi ed errori, e un protocollo di
valutazione del workflow, descritto separatamente nel Capitolo~\ref{cap:7}.

\section{Impostazione del progetto}

L'intera piattaforma (backend, frontend, server MCP, Redis, Keycloak e
harvester NebulaGraph) viene avviata come un insieme coordinato di container.
La domanda viene prima confrontata con il dominio geografico e ambientale, poi
trasformata in una rappresentazione strutturata e assegnata alle sorgenti
pertinenti. Un grafo LangGraph stabilisce i percorsi consentiti e permette ai
rami indipendenti di procedere in parallelo. La struttura rende ispezionabili
le condizioni applicate prima di raccogliere i risultati.

Ogni integrazione traduce la risposta remota in un dataset con campi stabili. I
client possono così essere provati con input controllati e gli errori restano
associati al servizio che li ha prodotti. Se un ramo fallisce, i risultati già
disponibili vengono conservati e la risposta segnala la fonte mancante.

La memoria è gestita come stato applicativo. A ogni conversazione è associato
un checkpoint Redis che comprende i messaggi necessari, i risultati delle
sorgenti e i riferimenti agli artefatti. Un turno successivo può così
riutilizzare un dataset o aggiornare soltanto la parte indicata dall'utente.
La memoria è propria di ciascuna conversazione e lascia invariati i parametri
del modello.

La verifica software copre parsing dei feed, filtri e aggregazioni con test
unitari; i percorsi del workflow vengono provati con dipendenze simulate; API,
persistenza e operazioni sulla dashboard sono controllate dai test di
integrazione. La valutazione end-to-end del comportamento dell'assistente è affidata al
protocollo~6.0 (dataset corrente da 105 domande, controlli deterministici e
LLM-as-a-Judge), descritto nel Capitolo~\ref{cap:7}. Dopo le fonti e
l'architettura, il Capitolo~\ref{cap:5} presenta i percorsi applicativi; i
Capitoli~\ref{cap:6} e~\ref{cap:7} approfondiscono rispettivamente
osservabilità e valutazione.

\section{Perimetro e limiti}

La descrizione dell'implementazione si basa sul codice del prototipo e sulla
guida tecnica del progetto \cite{guidaProgetto2026}.

Il sistema è un prototipo di ricerca: le comunicazioni ufficiali in caso di
emergenza restano di competenza delle autorità di protezione civile. GDACS
offre valutazioni rapide che possono cambiare quando vengono raccolte nuove
informazioni. Le attivazioni CEMS documentano prodotti richiesti e realizzati in
uno specifico contesto operativo. Il catalogo STAC indica quali osservazioni
sono disponibili; la stima dell'entità di un danno richiede analisi specifiche.

La dashboard descrive la disponibilità delle acquisizioni. La heatmap
rappresenta la densità dei centri delle bounding box degli item recuperati e i
rettangoli visualizzati mostrano le estensioni delle bounding box, che
approssimano le geometrie degli item. L'analisi scientifica delle immagini
(indici spettrali, classificazioni del suolo, stime degli effetti di un evento)
esula dal prototipo.

La copertura geografica dipende dai servizi consultati. Un luogo espresso in
linguaggio naturale deve essere risolto in coordinate o in una bounding box;
nomi ambigui possono richiedere un chiarimento. Ricerca web, cataloghi e feed
dipendono inoltre dalla rete e dalla disponibilità dei rispettivi fornitori: il
sistema gestisce errori e risultati parziali e segnala le fonti irraggiungibili.

Sul piano della valutazione, l'organizzazione multiagente è descritta e
verificata, ma manca ancora un confronto controllato con un agente singolo o
con una pipeline più semplice nelle stesse condizioni; un vantaggio
quantitativo di accuratezza, latenza o costo resta quindi da misurare. I
giudizi dell'LLM-as-a-Judge attendono una validazione umana e la verifica
end-to-end delle conversazioni nel browser è ancora da svolgere. Sul piano della
sicurezza, Keycloak protegge l'interfaccia, mentre le API del backend devono
ancora verificare il token e il proprietario degli artefatti.

Per questo le risposte conservano collegamenti e identificativi, distinguono i
dati recuperati dal testo generato e segnalano le sorgenti mancanti. Il criterio
adottato è rendere riconoscibile il percorso con cui il risultato è stato
costruito, anche quando alcune fonti sono assenti.
